\documentclass[11pt]{article}
\usepackage[T1]{fontenc}
\usepackage{lmodern,amsmath,amssymb,amsthm,mathtools}
\usepackage[margin=1in]{geometry}
\usepackage[colorlinks=true,linkcolor=blue,citecolor=blue,urlcolor=blue]{hyperref}
\newtheorem{theorem}{Theorem}
\newtheorem{lemma}{Lemma}
\newtheorem{corollary}{Corollary}
\newtheorem{example}{Example}
\newcommand{\ddc}{dd^c}
\newcommand{\cI}{\mathcal I}
\newcommand{\cO}{\mathcal O}
\newcommand{\PSH}{\operatorname{PSH}}
\newcommand{\ord}{\operatorname{ord}}
\newcommand{\PP}{\mathbb P}
\title{Fiber integrability and base multiplier-ideal approximation}
\author{}
\date{3 October 2026}
\begin{document}
\maketitle
\begin{abstract}
For a compact holomorphic submersion with connected fibers, we give sufficient
conditions under which sections selected by an upstairs singularity are
approximated by the multiplier ideals of one base potential of positive mass.
A global logarithmic lower bound and one anisotropic upper-bound chart place
the descended section spaces between divisorial series whose vanishing orders
differ by a bounded amount. Their dimension defect is $O(m^{q-1})$ on a base
of dimension $q$. Exact local models give the selected series explicitly,
with fiber-integrability gain $p\sum_j b_j^{-1}$. Product examples show that
the selecting singularity can be supported on a proper subset of each fiber.
The result answers the request for sufficient conditions in catalogue P21;
it gives no necessity-and-sufficiency classification for arbitrary selectors.
\end{abstract}

\section{The approximation question}
Let $\pi:X\to B$ be a surjective holomorphic submersion of connected compact
complex manifolds with connected fibers. Suppose that $B$ is K\"ahler, and
put $q=\dim B\geq1$ and $f=\dim X-q\geq1$. Fix a Hermitian form $\omega$
on $X$, a holomorphic line bundle $A$ with smooth Hermitian metric $h$ on
$B$, and $\varphi\in\PSH(X,\omega)$. Write $\theta=c_1(A,h)$.
We use $\ddc=\frac{i}{2\pi}\partial\bar\partial$ and define $\cI(tu)$ by
local integrability of $|g|^2e^{-tu}$ for holomorphic germs $g$.

The spaces in P21 of the attached catalogue~\cite{catalogue} are
\[
 V_m=H^0(X,\pi^*A^m\otimes\cI(m\varphi))\subset H^0(B,A^m).
\]
The identification follows from compact connected fibers. The question asks
for conditions producing $\psi\in\PSH(B,\theta)$ of positive
non-pluripolar mass such that
\[
 W_m=H^0(B,A^m\otimes\cI(m\psi))\subset V_m,
 \qquad \dim V_m-\dim W_m=o(m^q).
\]
No section norm is part of this assertion.

The isotropic section-over-divisor result is already proved in
Theorem 26.10.1 of Xia's working monograph~\cite{book}. Its error order and
sharp product example are established background. Here only one chart is
needed to detect the leading base vanishing order. Separate powers in that
chart measure the bounded relaxation supplied by fiber integration. A global
lower bound ensures that stronger selecting conditions elsewhere do not
destroy the base inclusion.

\section{A bounded divisorial band}
Let $D\subset B$ be a nonempty smooth connected reduced effective divisor.
Choose a smooth metric $H_D$ on $\cO(D)$, let $s_D$ be its defining section,
and write $\eta=c_1(\cO(D),H_D)$.

\begin{theorem}\label{thm:criterion}
Assume that there are $c>0$ and $u\in C^\infty(B,\mathbb R)$ such that
\begin{equation}\label{eq:positive}
 \beta=\theta+\ddc u-c\eta>0.
\end{equation}
Suppose that for some constant $C$ one has
\begin{equation}\label{eq:global}
 \varphi\geq c\log\|\pi^*s_D\|^2-C
 \quad\hbox{on }X,
\end{equation}
where the divisor norm uses the pulled-back metric $H_D$ and the inequality
is interpreted for extended-real values.

Suppose further that one coordinate polydisc on $X$ has submersion
coordinates
\[
 \pi(z,t,w,v)=(z,t),\quad D=\{z=0\},\quad
 t\in\mathbb C^{q-1},\ w\in\mathbb C^r,\ v\in\mathbb C^{f-r},
\]
with $1\leq r\leq f$, and throughout this polydisc
\begin{equation}\label{eq:witness}
 \varphi\leq\frac cp\log\left(|z|^{2p}+
                  \sum_{j=1}^r|w_j|^{2b_j}\right)+C'
\end{equation}
for positive integers $p,b_1,\ldots,b_r$. The chart lies over a nonempty
open subset of $D$.

Set
\[
 a=p\sum_{j=1}^r b_j^{-1},\qquad
 L_m=\lfloor mc\rfloor,\qquad
 K_m=\max\{\lfloor mc-a\rfloor,0\}.
\]
Then
\begin{equation}\label{eq:reference}
 \psi=u+c\log\|s_D\|^2\in\PSH(B,\theta),\qquad
 \int_B\langle(\theta+\ddc\psi)^q\rangle=\int_B\beta^q\in(0,\infty).
\end{equation}
For every positive integer $m$,
\begin{equation}\label{eq:band}
 W_m=H^0(B,A^m(-L_mD))\subset V_m
                  \subset H^0(B,A^m(-K_mD)),
\end{equation}
and
\begin{equation}\label{eq:defect}
 0\leq\dim V_m-\dim W_m=O(m^{q-1})=o(m^q).
\end{equation}
When $q=1$, the defect is at most $\lceil a\rceil$.
\end{theorem}

The positivity condition constructs the base reference. The two singularity
bounds then compare its ideal series with the selected series. The upper
bound needs neither a global section over $D$ nor a global normal-coordinate
splitting.

\subsection{Two elementary estimates}
\begin{lemma}\label{lem:integral}
Let $p,b_1,\ldots,b_r$ be positive integers, $P>0$, and $\ell\geq0$ an
integer. Near the origin in $\mathbb C^{r+1}$,
\[
 |z|^{2\ell}\left(|z|^{2p}+\sum_{j=1}^r|w_j|^{2b_j}\right)^{-P}
 \in L^1_{\rm loc}
 \quad\Longleftrightarrow\quad
 \frac{\ell+1}{p}+\sum_{j=1}^r\frac1{b_j}>P.
\]
Equality of the two sides in the last inequality gives logarithmic divergence.
\end{lemma}
\begin{proof}
Integrate the angular variables and put
$x_0=|z|^{2p}$ and $x_j=|w_j|^{2b_j}$. Up to positive constants the integral
becomes
\[
 \int (x_0+\cdots+x_r)^{-P}
 x_0^{(\ell+1)/p-1}\prod_{j=1}^r x_j^{1/b_j-1}\,dx_0\cdots dx_r.
\]
A small simplex $\sum x_j<\delta$ is contained in the integration box, and
the box is contained in a larger simplex. On either simplex use
$\rho=\sum x_j$ and $x_j=\rho y_j$. The angular simplex integral is finite
and positive since all its exponents are greater than $-1$. The radial
integral is
\[
 \int_0^\delta\rho^{(\ell+1)/p+\sum b_j^{-1}-P-1}\,d\rho.
\]
This proves the assertion and the equality endpoint. Additional tangential
variables on bounded polydiscs, bounded remainders in the weight, and a
holomorphic unit multiplying $z^\ell$ do not change local integrability.
\end{proof}

\begin{lemma}\label{lem:dimension}
Let $Y$ be a compact complex manifold of dimension $s$ and let $E,F$ be
fixed holomorphic line bundles. Uniformly for integers $j,k$,
\[
 h^0(Y,E^j\otimes F^k)=O((1+|j|+|k|)^s).
\]
Consequently, if $0\leq K_m\leq L_m=O(m)$ and $L_m-K_m$ is bounded, then
\[
 h^0(B,A^m(-K_mD))-h^0(B,A^m(-L_mD))=O(m^{q-1}).
\]
\end{lemma}
\begin{proof}
Choose smooth bundle metrics, a smooth positive volume, and finitely many
coordinate charts and frames with smaller charts covering $Y$. Put
$N=1+|j|+|k|$. All second derivatives of the metric weight of
$E^j\otimes F^k$ are $O(N)$. At a point in a smaller chart, a holomorphic
affine exponential change of frame removes the constant and real-linear
terms of that weight. The new weight is uniformly bounded on a ball of
radius $\epsilon N^{-1/2}$ about the point. Submean therefore gives
\[
 |\sigma(y)|^2\leq C N^s\|\sigma\|_{L^2}^2.
\]
For an orthonormal basis, the sum of the squared evaluations is the squared
operator norm of evaluation into a one-dimensional fiber. It is thus at
most $CN^s$. Integrating the sum bounds the dimension. In dimension zero
the assertion follows from the finite number of points.

For the second assertion use the restriction sequences
\[
 0\longrightarrow A^m(-(j+1)D)\longrightarrow A^m(-jD)
 \longrightarrow A^m(-jD)|_D\longrightarrow0
\]
for $K_m\leq j<L_m$. Each successive dimension difference is at most
$h^0(D,A^m(-jD)|_D)$. Apply the first assertion on the compact complex
$(q-1)$-fold $D$ with $j=O(m)$, and sum over the bounded number of steps.
No surjectivity of restriction or cohomology vanishing is required. If
$q=1$, the connected smooth divisor is one point and each restriction space
has dimension one.
\end{proof}

\subsection{Proof of the criterion}
\begin{proof}[Proof of Theorem \ref{thm:criterion}]
In a local frame of $A$, a section of $\pi^*A^m$ is a holomorphic function
on each compact connected fiber, and is consequently constant there. Local
holomorphic right inverses of the submersion show that its descended
coefficient is holomorphic. These coefficients glue, giving
\[
 H^0(X,\pi^*A^m)=H^0(B,A^m).
\]

Poincar\'e--Lelong gives
$\theta+\ddc\psi=\beta+c[D]\geq0$. Off $D$ the current is $\beta$.
The non-pluripolar construction discards the locus $\psi=-\infty$, which
is $D$, so its top product is the extension of $\beta^q$ from $B\setminus D$.
Since a smooth volume gives $D$ zero measure, that extension is $\beta^q$.
Strict positivity and compactness prove \eqref{eq:reference}, including
when $q=1$.

Locally $\psi=c\log|z|^2+O(1)$. For a germ of divisor order $\ell$,
one-variable integrability requires $\ell+1>mc$. Thus
\[
 \cI(m\psi)=\cO(-\lfloor mc\rfloor D).
\]
In particular, if $mc$ is an integer, order $mc-1$ fails and order $mc$
succeeds. The bounded function $u$ and \eqref{eq:global} give
$\varphi\geq\pi^*\psi-C''$. In submersion coordinates, integrability
on the base implies integrability of the pullback on a bounded fiber
polydisc. Therefore $W_m\subset V_m$ for every $m$.

Take a nonzero descended section $\sigma\in V_m$ and put
$\ell=\ord_D(\sigma)$. Since the smooth connected divisor is irreducible,
in the witness chart its coefficient is $z^\ell g(z,t)$ with $g(0,t)$
not identically zero. Otherwise the identity theorem on $D$ would increase
its divisor order. Choose a smaller chart where $g$ is a unit.
The upper bound \eqref{eq:witness} implies that integrability of
$|\pi^*\sigma|^2e^{-m\varphi}$ requires integrability of
\[
 |z|^{2\ell}\left(|z|^{2p}+\sum|w_j|^{2b_j}\right)^{-mc/p}.
\]
Lemma \ref{lem:integral} gives $\ell+1>mc-a$. For integer $\ell\geq0$
this is exactly $\ell\geq\max\{\lfloor mc-a\rfloor,0\}=K_m$.
When $q=1$ the coefficient $g$ has no $t$ variables and its leading value
is nonzero. This proves the upper inclusion of \eqref{eq:band}.

Finally
\[
 0\leq L_m-K_m\leq\lceil a\rceil,
\]
by the elementary floor inequality and truncation at zero. Both orders
are $O(m)$. Lemma \ref{lem:dimension} and \eqref{eq:band} give
\eqref{eq:defect}. Its dimension-zero restriction argument also proves the
stated bound when $q=1$.
\end{proof}

\section{Exact models and non-descended examples}
\begin{corollary}\label{cor:exact}
Retain \eqref{eq:positive}. Suppose $Z\subset\pi^{-1}D$ is a nonempty
compact smooth submanifold projecting submersively onto $D$, and that near
each point of $Z$ there are the submersion coordinates above with
\[
 Z=\{z=w_1=\cdots=w_r=0\},\qquad
 \varphi=\frac cp\log\left(|z|^{2p}+\sum_{j=1}^r|w_j|^{2b_j}\right)+O(1),
\]
with fixed $p,b_j,c$. Suppose $\varphi$ is locally bounded outside $Z$.
Then the conclusions of Theorem \ref{thm:criterion} hold and, for every $m$,
\[
 V_m=H^0(B,A^m(-K_mD)).
\]
\end{corollary}
\begin{proof}
The local model dominates $c\log|z|^2$ up to a constant. Away from $Z$,
local boundedness gives the same lower bound near $\pi^{-1}D$, and outside
$\pi^{-1}D$ both functions are locally bounded. A finite covering of
compact $X$ proves \eqref{eq:global}. The upper witness is immediate.

For the remaining inclusion, a section divisible by $z^{K_m}$ has a bounded
holomorphic residual coefficient in every pole chart. The strict inequality
$K_m+1>mc-a$ and Lemma \ref{lem:integral} imply local integrability there.
Outside $Z$ the weight is locally bounded. Thus every section of
$H^0(B,A^m(-K_mD))$ belongs to $V_m$.
\end{proof}

For a nonempty family, take $X=B\times(\PP^1)^f$ with a projective base,
an ample $A$, and the divisor $D$ above. In the first $r$ fiber factors mark
points $o_j$ and let $E_j$ be their inverse-image divisors. Choose smooth
metrics on the corresponding divisor bundles and set
\[
 Q=\|\pi^*s_D\|^{2p}+\sum_{j=1}^r\|s_{E_j}\|^{2b_j},\qquad
 \varphi=\lambda\log Q,\qquad c=p\lambda.
\]
The common zero set is
$D\times\{o_1\}\times\cdots\times\{o_r\}\times(\PP^1)^{f-r}$.
Near it, smooth positive metric factors give the exact model of
Corollary \ref{cor:exact}; elsewhere the function is smooth.
The inequality $Q\geq\|\pi^*s_D\|^{2p}$ gives \eqref{eq:global}.

There is a constant $C_0$ with $\ddc\log Q\geq-C_0\omega$. Indeed, $Q$
is the squared norm of a holomorphic section of the direct sum of the
divisor line bundles with powers $p,b_j$. Locally the log-sum-exp Hessian
is a weighted average of their smooth negative curvatures plus a positive
variance term. A local quadratic correction gives a psh function away from
the component zeros, and the locally bounded-above psh extension preserves
the inequality across the common zero set. Compactness makes $C_0$ uniform.
Thus $\lambda C_0\leq1$ ensures $\varphi\in\PSH(X,\omega)$ for the
fixed Hermitian background. A positive metric on $A$ differs smoothly from
the prescribed $h$. Choose $u$ giving its curvature; for sufficiently small
$c=p\lambda$, \eqref{eq:positive} holds as well. Both smallness conditions
can be imposed simultaneously.

These potentials are singular on a proper subset of each fiber over $D$ and
finite elsewhere on that fiber. Hence they cannot be pullbacks of a base
potential up to bounded difference.

\begin{example}\label{ex:explicit}
Let $B=\PP^q$, let $D$ be a hyperplane, take $A=\cO(1)$, and put
$X=B\times\PP^1$. Use Fubini--Study divisor metrics and choose a point
$o\in\PP^1$. Set
\[
 \varphi=\tfrac14\log\bigl(\|s_D\|^4+\|s_o\|^6\bigr),\qquad
 p=2,\quad b_1=3,\quad c=\tfrac12.
\]
For $\omega=\omega_B+\omega_F$, the curvature estimate just used gives
$\ddc\varphi\geq-\frac12\omega_B-\frac34\omega_F$, so the selector
is admissible. The residual is $\beta=\frac12\omega_B$, with mass $2^{-q}$
when $\int_B\omega_B^q=1$.

Here $a=2/3$, $L_m=\lfloor m/2\rfloor$, and
$K_m=\max\{\lfloor m/2-2/3\rfloor,0\}$. For $m=1$ both selected
spaces are $H^0(\cO(1))$. For every $m\geq2$, $K_m=L_m-1$, and
\[
 V_m=H^0(\PP^q,\cO(m-L_m+1)),\qquad
 W_m=H^0(\PP^q,\cO(m-L_m)).
\]
Consequently
\[
 \dim V_m-\dim W_m
 =\binom{\lceil m/2\rceil+q}{q-1}
 \sim\frac{2^{1-q}}{(q-1)!}m^{q-1}.
\]
For $q=1$ the defect is exactly one for $m\geq2$. Thus the error order
is sharp in this anisotropic class.
\end{example}

The common horizontal coefficient $c$, a witness over an open part of $D$,
and the finite weighted transverse gain $p\sum b_j^{-1}$ suffice for the
approximation. The global lower bound excludes additional constraints that
would obstruct the inclusion. With $p=b_j=1$ and $r=f$, the exact order is
$\max\{\lfloor mc\rfloor-f,0\}$, recovering the established isotropic
section-over-divisor case. Separate powers, positive-dimensional pole sets
inside the fibers, and the one-chart sandwich criterion extend that regime.

If the fibers have dimension zero, the connected-fiber submersion is a
biholomorphism. When the descended selector itself is $\theta$-psh of finite
positive mass, taking it as $\psi$ gives $V_m=W_m$. The examples above concern
the additional integrability supplied by positive-dimensional fibers.

The positivity assumption cannot simply be omitted. If $A$ is trivial with
flat metric, global plurisubharmonic functions on compact connected $B$
are constant and have zero non-pluripolar top mass. The theorem provides
explicit sufficient conditions for P21 and leaves arbitrary singularities,
more general base images, and necessity questions open.

\begin{thebibliography}{9}
\bibitem{catalogue}
\emph{Open problems in non-K\"ahler pluripotential theory: full statements,
known results, and remaining gaps}, catalogue accompanying Mingchen Xia's
monograph, 2 October 2026, supplied as \texttt{document.pdf}, P21.
\bibitem{book}
M. Xia, \emph{Singularities in Hermitian pluripotential theory}, public
working manuscript, \url{https://mingchenxia.github.io/Papers/SHPT_vibe.pdf},
Lemma 26.10.1, Theorem 26.10.1, and Example 26.10.1.
\end{thebibliography}
\end{document}
